\section{Going down for integral over normal}
\label{section-going-down-integral-over-normal}

\noindent
We first play around a little bit with the notion of elements
integral over an ideal, and then we prove the theorem referred
to in the section title.

\begin{definition}
\label{definition-integral-over-ideal}
Let $\varphi : R \to S$ be a ring map.
Let $I \subset R$ be an ideal.
We say an element $g \in S$ is
{\it integral over $I$} if
there exists a monic
polynomial $P = x^d + \sum_{j < d} a_j x^j$
with coefficients $a_j \in I^{d-j}$ such
that $P^\varphi(g) = 0$ in $S$.
\end{definition}

\noindent
This is mostly used when $\varphi = \text{id}_R : R \to R$.
In this case the set $I'$ of elements integral over $I$ is called
the {\it integral closure of $I$}. We will see that $I'$ is
an ideal of $R$ (and of course $I \subset I'$).

\begin{lemma}
\label{lemma-characterize-integral-ideal}
Let $\varphi : R \to S$ be a ring map.
Let $I \subset R$ be an ideal.
Let $A = \sum I^nt^n \subset R[t]$ be the
subring of the polynomial ring
generated by $R \oplus It \subset R[t]$.
An element $s \in S$ is integral over $I$ if
and only if the element $st \in S[t]$
is integral over $A$.
\end{lemma}

\begin{proof}
Suppose $st$ is integral over $A$.
Let $P = x^d + \sum_{j < d} a_j x^j$
be a monic polynomial with coefficients in $A$
such that $P^\varphi(st) = 0$. Let $a_j' \in A$
be the degree $d-j$ part of $a_j$, in other
words $a_j' = a_j'' t^{d-j}$ with $a_j'' \in I^{d-j}$.
For degree reasons we still have
$(st)^d + \sum_{j < d} \varphi(a_j'') t^{d-j} (st)^j = 0$.
Hence $s^d + \sum_{j < d} \varphi(a_j'') s^j = 0$
and we see that $s$ is integral over $I$.

\medskip\noindent
Suppose that $s$ is integral over $I$.
Say $P = x^d + \sum_{j < d} a_j x^j$
with $a_j \in I^{d-j}$. Then we immediately find a
polynomial $Q = x^d + \sum_{j < d} (a_j t^{d-j}) x^j$
with coefficients in $A$ which proves that
$st$ is integral over $A$.
\end{proof}

\begin{lemma}
\label{lemma-integral-over-ideal-is-submodule}
Let $\varphi : R \to S$ be a ring map.
Let $I \subset R$ be an ideal.
The set of elements of $S$ which are integral
over $I$ form a $R$-submodule of $S$.
Furthermore, if $s \in S$ is integral over
$R$, and $s'$ is integral over $I$, then
$ss'$ is integral over $I$.
\end{lemma}

\begin{proof}
We will use Lemma \ref{lemma-integral-closure-is-ring} without
further mention. Closure under addition is clear from the
characterization of Lemma \ref{lemma-characterize-integral-ideal}
whose notation we adopt.
Any element $s \in S$ which is integral over
$R$ corresponds to the degree $0$ element $s$ of $S[t]$
which is integral over $A$ (because $R \subset A$).
Hence we see that multiplication by $s$ on $S[t]$
preserves the property of being integral over $A$,
\end{proof}

\begin{lemma}
\label{lemma-integral-integral-over-ideal}
Suppose $\varphi : R \to S$ is integral.
Suppose $I \subset R$ is an ideal.
Then every element of $IS$ is integral over $I$.
\end{lemma}

\begin{proof}
Immediate from Lemma \ref{lemma-integral-over-ideal-is-submodule}.
\end{proof}

\begin{lemma}
\label{lemma-polynomials-divide}
Let $K$ be a field. Let $n, m \in \mathbf{N}$ and
$a_0, \ldots, a_{n - 1}, b_0, \ldots, b_{m - 1} \in K$.
If the polynomial $x^n + a_{n - 1}x^{n - 1} + \ldots + a_0$
divides the polynomial $x^m + b_{m - 1} x^{m - 1} + \ldots + b_0$
in $K[x]$ then
\begin{enumerate}
\item $a_0, \ldots, a_{n - 1}$ are integral over any subring
$R_0$ of $K$ containing the elements $b_0, \ldots, b_{m - 1}$, and
\item each $a_i$ lies in $\sqrt{(b_0, \ldots, b_{m-1})R}$
for any subring $R \subset K$ containing the elements
$a_0, \ldots, a_{n - 1}, b_0, \ldots, b_{m - 1}$.
\end{enumerate}
\end{lemma}

\begin{proof}
Let $L/K$ be a field extension such that we can write
$x^m + b_{m - 1} x^{m - 1} + \ldots + b_0 =
\prod_{i = 1}^m (x - \beta_i)$ with $\beta_i \in L$.
See Fields, Section \ref{fields-section-splitting-fieds}.
Each $\beta_i$ is integral over $R_0$.
Since each $a_i$ is a homogeneous polynomial in $\beta_1, \ldots, \beta_m$
we deduce the same for the $a_i$
(use Lemma \ref{lemma-integral-closure-is-ring}). This proves (1).

\medskip\noindent
Let $R$ be as in (2). Choose $c_0, \ldots, c_{m - n - 1} \in K$ such that
$$
\begin{matrix}
x^m + b_{m - 1} x^{m - 1} + \ldots + b_0 =  \\
(x^n + a_{n - 1}x^{n - 1} + \ldots + a_0)
(x^{m - n} + c_{m - n - 1}x^{m - n - 1}+ \ldots + c_0).
\end{matrix}
$$
This equation implies
\begin{align*}
c_{m - n - 1} & = b_{m - 1} - a_{n - 1}, \\
c_{m - n - 2} & = b_{m - 2} - a_{n - 2} - a_{n - 1}c_{m - n - 1}, \\
\ldots
\end{align*}
Thus $c_j \in R$ for all $j$. Dividing out the radical
$\sqrt{(b_0, \ldots, b_{m - 1})}$ we get a reduced ring $\overline{R}$.
We have to show that the images $\overline{a}_i \in \overline{R}$
are zero. And in
$\overline{R}[x]$ we have the relation
$$
\begin{matrix}
x^m = x^m + \overline{b}_{m - 1} x^{m - 1} + \ldots + \overline{b}_0 = \\
(x^n + \overline{a}_{n - 1}x^{n - 1} + \ldots + \overline{a}_0)
(x^{m - n} + \overline{c}_{m - n - 1}x^{m - n - 1}+ \ldots + \overline{c}_0).
\end{matrix}
$$
It is easy to see that this implies $\overline{a}_i = 0$ for all $i$. Indeed
by Lemma \ref{lemma-minimal-prime-reduced-ring} the localization of
$\overline{R}$ at a minimal prime $\mathfrak{p}$ is a field and
$\overline{R}_{\mathfrak p}[x]$ a UFD. Thus
$f = x^n + \sum \overline{a}_i x^i$
is associated to $x^n$ and since $f$ is monic $f = x^n$
in $\overline{R}_{\mathfrak p}[x]$.
Then there exists an $s \in \overline{R}$, $s \not\in \mathfrak p$
such that $s(f - x^n) = 0$.  Therefore all $\overline{a}_i$ lie
in $\mathfrak p$ and we conclude by
Lemma \ref{lemma-reduced-ring-sub-product-fields}.
\end{proof}

\begin{lemma}
\label{lemma-minimal-polynomial-normal-domain}
Let $R \subset S$ be an inclusion of domains.
Assume $R$ is normal. Let $g \in S$ be integral
over $R$. Then the minimal polynomial of $g$
has coefficients in $R$.
\end{lemma}

\begin{proof}
Let $P = x^m + b_{m-1} x^{m-1} + \ldots + b_0$
be a polynomial with coefficients in $R$
such that $P(g) = 0$. Let $Q = x^n + a_{n-1}x^{n-1} + \ldots + a_0$
be the minimal polynomial for $g$ over the fraction field
$K$ of $R$. Then $Q$ divides $P$ in $K[x]$. By Lemma
\ref{lemma-polynomials-divide} we see the $a_i$ are
integral over $R$. Since $R$ is normal this
means they are in $R$.
\end{proof}

\begin{proposition}
\label{proposition-going-down-normal-integral}
Let $R \subset S$ be an inclusion of domains.
Assume $R$ is normal and $S$ integral over $R$.
Let $\mathfrak p \subset \mathfrak p' \subset R$
be primes. Let $\mathfrak q'$ be a prime of $S$
with $\mathfrak p' = R \cap \mathfrak q'$.
Then there exists a prime $\mathfrak q$
with $\mathfrak q \subset \mathfrak q'$
such that $\mathfrak p = R \cap \mathfrak q$. In other words:
the going down property holds for $R \to S$, see
Definition \ref{definition-going-up-down}.
\end{proposition}

\begin{proof}
Let $\mathfrak p$, $\mathfrak p'$ and $\mathfrak q'$
be as in the statement. We have to show there is a prime
$\mathfrak q$, with $\mathfrak q \subset \mathfrak q'$ and
$R \cap \mathfrak q = \mathfrak p$. This is the same
as finding a prime of
$S_{\mathfrak q'}$ mapping to $\mathfrak p$.
According to Lemma \ref{lemma-in-image} we have to show
that $\mathfrak p S_{\mathfrak q'} \cap R
= \mathfrak p$. Pick $z \in \mathfrak p S_{\mathfrak q'} \cap R$.
We may write $z = y/g$ with $y \in \mathfrak pS$ and
$g \in S$, $g \not\in \mathfrak q'$. Written
differently we have $zg = y$.

\medskip\noindent
By Lemma \ref{lemma-integral-integral-over-ideal}
there exists a monic polynomial
$P = x^m + b_{m-1} x^{m-1} + \ldots + b_0$
with $b_i \in \mathfrak p$ such that $P(y) = 0$.

\medskip\noindent
By Lemma \ref{lemma-minimal-polynomial-normal-domain}
the minimal polynomial of $g$ over $K$ has coefficients
in $R$. Write it as $Q = x^n + a_{n-1} x^{n-1} + \ldots
+ a_0$. Note that not all $a_i$, $i = n-1, \ldots, 0$
are in $\mathfrak p$ since that would imply
$g^n = \sum_{j < n} a_j g^j \in \mathfrak pS
\subset \mathfrak p'S
\subset \mathfrak q'$
which is a contradiction.

\medskip\noindent
Since $y = zg$ we see immediately from the above
that $Q' = x^n + za_{n-1} x^{n-1} + \ldots + z^{n}a_0$
is the minimal polynomial for $y$. Hence
$Q'$ divides $P$ and by Lemma \ref{lemma-polynomials-divide}
we see that $z^ja_{n - j} \in \sqrt{(b_0, \ldots, b_{m-1})}
\subset \mathfrak p$, $j =  1, \ldots, n$.
Because not all $a_i$, $i = n-1, \ldots, 0$
are in $\mathfrak p$ we conclude $z \in \mathfrak p$
as desired.
\end{proof}














